\documentclass[../../../include/open-logic-section]{subfiles}

\begin{document}

\olfileid{sth}{choice}{countablechoice}
\olsection{可数选择公理}

在不使用任何形式的选择公理/良序原理的情况下，很容易证明：

\begin{lem}[在 $\Zminus$ 中]
每个有限集都有一个选择函数。
\end{lem}

\begin{proof}
设 $a = \{b_1, \ldots, b_n\}$。为简化起见，假设每个 $b_i
\neq \emptyset$。因此存在对象 $c_1, \ldots, c_n$ 使得
$c_1 \in b_1, \ldots, c_n \in b_n$。现在根据
\olref[z][pairs]{prop:pairsconsequences}，集合 $\{\langle b_1,
c_1\rangle , \ldots, \langle b_n, c_n\rangle\}$ 存在；这就是~$a$ 的一个
选择函数。
\end{proof}

但一旦我们考虑无限集，情况就更不明朗了。例如，考虑对上述内容的以下“最小”扩展：

\begin{defish}
\emph{可数选择公理。}每个\emph{可数}集都有一个选择函数。
\end{defish}

这是选择公理的一个特例。事实表明，这一原理过去曾被相当频繁地使用，但使用者往往并未明确意识到自己使用了它。这里有两个很好的例子。\footnote{分别源于
\citet[\S9.4]{Potter2004} 和 Luca Incurvati。}

\begin{ex}
这里有一个自然的想法：对任意集合 $A$，要么
$\cardle{\omega}{A}$，要么对某个 $n \in \omega$ 有 $\cardeq{A}{n}$。这是
表达“每个集合要么有限要么无限”这一直观想法的一种方式。Cantor 和许多其他数学家都曾断言这一点而未加证明。尽管我们很谨慎，在
\olref[cardinals][classing]{generalinfinitycharacter} 中已经证明了它。
但在那个证明中，我们在 $\ZFC$ 中进行工作，因为我们假定任何集合 $A$ 都可以良序，从而保证了 $\card{A}$ 存在。也就是说：我们明确地假定了选择公理。

事实上，\citet{Dedekind1888} 为这一断言提供了他自己的证明，如下：

\begin{thm}[在 $\Zminus + \text{可数选择公理}$ 中]
对任意 $A$，要么 $\cardle{\omega}{A}$，要么对某个 $n \in \omega$ 有 $\cardeq{A}{n}$。
\end{thm}

\begin{proof}
假设对所有 $n \in \omega$ 都有 $\cardneq{A}{n}$。那么特别地，对每个
$n < \omega$ 存在子集 $A_n \subseteq A$，恰好有 $\cardexpo{2}{n}$ 个元素。利用这个序列 $A_0, A_1, A_2,
\ldots$，我们对每个 $n$ 定义：
\[
	B_n = A_n \setminus \bigcup_{i < n} A_i.
\]
现在注意到
\begin{align*}
	\card{\bigcup_{i < n}A_n}
	&\leq \card{A_0} + \card{A_1} + \ldots + \card{A_{n-1}}\\
	&=1 + 2 + \ldots + 2^{n-1}\\
	& = 2^n - 1\\
	& < 2^n = \card{A_n}
\end{align*}
因此每个 $B_n$ 都至少有一个成员 $c_n$。此外，这些 $B_n$ 两两不交；所以如果 $c_n = c_m$ 则 $n = m$。且每个 $c_n
\in A$。所以函数  $f(n) = c_n$ 是一个从 $\omega$ 到 $A$ 的单射。
\end{proof}
\noindent
Dedekind 没有明确指出他使用了可数选择公理。但是，\emph{你}发现它的使用了吗？再看一遍。（真的：\emph{再看一遍}。）

这个证明两次使用了可数选择公理。我们第一次使用它，是为了得到集合序列 $A_0$, $A_1$, $A_2$, \dots\@ 我们第二次使用它，是为了从每个~$B_n$ 中选出元素 $c_n$。而且，这里对选择公理的使用是必不可少的。\citet[p.~138]{Cohen1966} 证明了，如果我们没有任何形式的选择公理，这个结论就不成立。也就是说：存在与 $\ZF$ 一致的模型，其中包含与~$\omega$ 不可比较的集合。
\end{ex}

\begin{ex}
在 \citeyear{Cantor1878}，Cantor 声称可数多个可数集合的并是可数的。他并未给出证明，或许是因为他认为这个证明是显然的。尽管我们很谨慎，在
\olref[card-arithmetic][simp]{kappaunionkappasize} 中证明了这个结果的一个更一般的形式。但我们的证明明确使用了选择公理。甚至这个较弱结果的证明也需要可数选择公理。

\begin{thm}[在 $\Zminus + \text{可数选择公理}$ 中]
如果对每个 $n \in \omega$，$A_n$ 都是可数的，那么 $\bigcup_{n <
\omega} A_n$ 是可数的。
\end{thm}

\begin{proof}
不失一般性，假设每个 $A_n \neq \emptyset$。因此对每个 $n \in \omega$ 存在一个满射 $f_n \colon \omega
\to A_n$。定义 $f \colon \omega \times \omega \to \bigcup_{n <
\omega} A_n$ 为 $f(m, n) = f_n(m)$。因为 $\omega
\times \omega$ 是可数的
（\olref[sfr][siz][zigzag]{natsquaredenumerable}）且 $f$ 是满射，结论成立。
\end{proof}
\noindent
你发现可数选择公理的使用了吗？它被用来选择函数序列 $f_0$, $f_1$, $f_2$, \dots\footnote{类似地使用选择公理的情况出现在
\olref[card-arithmetic][simp]{kappaunionkappasize}，当时我们指示“对每个 $\beta \in \cardfont{a}$，固定一个单射
$f_\beta$”。} 而且在没有任何选择公理的情况下，此结论不成立。具体而言，\citet{FefermanLevy1963} 证明了在 $\ZF$ 中，可数多个可数集合的并具有势~$\beth_1$ 是相容的。但关于这一点，有一个有趣得多的说法，来自 Russell：
\begin{quote}
  这可以用一个百万富翁的例子来说明：他每次买一双靴子时都会买一双袜子，而且只在那个时候买，他对购买两者如此热衷，以至于最后他有
  $\aleph_0$ 双靴子和 $\aleph_0$ 双袜子\dots\@ 在靴子中，我们可以区分左和右，因此我们可以从每双中选出一只，即我们可以选择所有右靴或所有左靴；但对于袜子，没有这样的选择原则会自然呈现，而且除非我们假设乘法公理 [即，实际上就是选择公理]，否则我们不能确定是否存在一个类，它由每双袜子中的一只组成。
  \citep[p.~126]{Russell1919}
\end{quote}
简而言之，需要某种形式的选择公理来证明以下结论：如果你有可数多双袜子，那么你（只）有可数多只袜子。事实上，没有可数选择公理（或等价命题），可数多个可数集合的并可能不是可数的。
\end{ex}

结论是，可数选择公理被反复使用，但其使用者大多没有意识到。哲学问题是：我们如何能\emph{为}可数选择公理\emph{提供辩护}？

一种直觉上的辩护尝试可能会诉诸于超任务。假设我们在 $\nicefrac{1}{2}$ 分钟内做出第一个选择，在 $\nicefrac{1}{4}$ 分钟内做出第二个选择，\dots，在第 $n$ 次选择时花费 $\nicefrac{1}{2^n}$ 分钟，\dots\@ 那么在 $1$~分钟内，我们将完成一个 $\omega$ 序列的选择，并定义出一个选择函数。

但是，这样的思想实验究竟能告诉我们什么？首先，它依赖于过于字面地理解“选择”这一概念。其次，它似乎将数学与形而上学的可能性绑定在一起。

更重要的是：它并不能为我们提供选择公理本身（而\emph{不仅仅}是可数选择公理）的辩护。因为如果我们需要\emph{每个}集合都有选择函数，那么我们将需要能够执行“任意序数长度的超任务”。坦率地说，这个想法是荒谬的。

\end{document}